\section{Couniversality and the unit on posets}\label{sec:couniversality}

Fix $n\in\{1,2,\ldots\}\cup\{\infty\}$ throughout this section.  A
\emph{Thomason equivalence} is a strict functor whose Street nerve is a
weak homotopy equivalence.  This terminology uses no model structure on
$\nCat$.  We prove that certain poset inclusions remain Thomason
equivalences after every pushout, and then establish the adjunction unit
on posets independently of transfer.

For an inclusion $i:A\hookrightarrow E$ of posets, a right adjoint
retraction $r:E\to A$ satisfies
\begin{equation}\label{eq:poset-adjunction}
  a\le e\quad\Longleftrightarrow\quad a\le r(e)
  \qquad(a\in A,\ e\in E),
  \qquad ri=\id_A.
\end{equation}
We identify $A$ with its image in $E$.  In particular $r(e)\le e$.
The inclusion is a \emph{sieve} when $e\le a\in A$ implies $e\in A$.
For an arbitrary strict functor $a:\U{A}\to X$, write
\begin{equation}\label{eq:couniversal-pushout}
  Y=X\amalg_{\U{A}}\U{E},\qquad
  j:X\longrightarrow Y,\qquad v:\U{E}\longrightarrow Y.
\end{equation}
The symbol $a$ here denotes the attaching functor; elements of $A$ will
be denoted by $b$ below.

\begin{lemma}[A greatest new element]\label{lem:greatest-new}
Suppose that $A\hookrightarrow E$ is a sieve with a right adjoint
retraction, and that $E\setminus A$ has a greatest element.  For every
attaching functor in \eqref{eq:couniversal-pushout}, the map $j$ admits a
retraction $p:Y\to X$ such that $\Street(jp)$ and $\id_{\Street Y}$ are
joined by a zigzag of simplicial homotopies.  Consequently $j$ is a
Thomason equivalence.  Neither $A$ nor $E$ needs to be finite.
\end{lemma}

\begin{proof}
Let $x$ be greatest in $E\setminus A$ and put $t=r(x)\in A$.
If $b\in A$, $y\in E\setminus A$, and $b\le y$, then
$b\le y\le x$, so \eqref{eq:poset-adjunction} gives $b\le t$.
Define maps $k,g:E\to E$ by
\[
 k(e)=\begin{cases}e,&e\in A,\\t,&e\notin A,\end{cases}
 \qquad
 g(e)=\begin{cases}e,&e\in A,\\x,&e\notin A.\end{cases}
\]
Both are monotone.  Indeed, for a comparable pair $e\le e'$, the
old--old case is unchanged and the new--new case is constant.  In the
old--new case the preceding observation gives $e\le t\le x$.
The new--old case cannot occur because $A$ is a sieve.  We therefore
have pointwise inequalities
\begin{equation}\label{eq:local-zigzag}
  \id_E\le g,\qquad k\le g,
\end{equation}
and all three maps restrict to $\id_A$.  The common upper map $g$ is
constant on $E\setminus A$, exactly as required in
Lemma~\ref{lem:relative-prism}.
Figure~\ref{fig:local-diamond} shows why this common upper map is useful:
$\id_E$ and $k$ need not be pointwise comparable, although both lie
below $g$.

Apply Lemma~\ref{lem:pushout-cylinder} to the two comparisons in
\eqref{eq:local-zigzag}.  Let $G,K:Y\to Y$ be the pushout endomorphisms
whose restrictions to $X$ are $j$ and whose restrictions to $\U{E}$
are $v\,c_n\OrdN(g)$ and $v\,c_n\OrdN(k)$, respectively.  The resulting
right Gray cylinders, both stationary on $X$, have endpoint pattern
\[
  \id_Y\ \longrightarrow\ G\ \longleftarrow\ K.
\]
By Lemma~\ref{lem:gray-nerve}, these give simplicial homotopies with the
same endpoint pattern on Street nerves.

The image of $k$ lies in $A$, so $k=i\bar k$ for a monotone retraction
$\bar k:E\to A$.  The maps $\id_X$ and
$a\,c_n\OrdN(\bar k):\U{E}\to X$ agree on $\U{A}$ and hence define
$p:Y\to X$.  They give $pj=\id_X$.  Comparing restrictions to the two
pushout summands gives $jp=K$: on $X$ both restrict to $j$, and on
$\U{E}$ one has
\[
 jp\,v=j\,a\,c_n\OrdN(\bar k)
       =v\,c_n\OrdN(i)\,c_n\OrdN(\bar k)
       =v\,c_n\OrdN(k).
\]
After realization, concatenate the homotopy from $\id_Y$ to $G$ with
the reverse of the homotopy from $jp$ to $G$.  Thus $|\Street p|$ is
a homotopy inverse to $|\Street j|$, proving the assertion.  This
reversal is a reversal of an ordinary topological homotopy; no reversed
Gray cylinder is asserted.
\end{proof}

\begin{figure}[htbp]
\centering
\begin{tikzpicture}[>=Stealth,scale=.9,
 old/.style={draw=blue!65!black,fill=blue!10,rounded corners=2pt,
             minimum width=7mm,minimum height=7mm},
 new/.style={draw,circle,minimum size=7mm}]
 \node[old] (a) at (0,0) {$a$};
 \node[old] (b) at (-1.2,1.1) {$b$};
 \node[new] (y) at (1.2,1.1) {$y$};
 \node[new] (x) at (0,2.2) {$x$};
 \draw[->] (a)--(b);
 \draw[->] (a)--(y);
 \draw[->] (b)--(x);
 \draw[->] (y)--(x);
 \node at (5.4,1.1)
   {$\id_E(y)=y\;\le\;g(y)=x\;\ge\;k(y)=b$};
\end{tikzpicture}
\caption{The shaded square nodes form the old sieve $A=\{a,b\}$.
Its right adjoint retraction has $r(y)=a$ and $r(x)=t=b$.
Thus $k(y)=b$ is incomparable with $y$; the common upper value $x$
supplies the two comparisons used by the relative prisms.}
\label{fig:local-diamond}
\end{figure}

\begin{theorem}[Couniversality]\label{thm:couniversal}
If $i:A\hookrightarrow E$ is a sieve of small posets with a right
adjoint retraction, then $c_n\OrdN(i)$ is a couniversal Thomason
equivalence: for every $a:\U{A}\to X$, the canonical map
\[
 X\longrightarrow X\amalg_{\U{A}}\U{E}
\]
is a Thomason equivalence.  This holds for every finite $n\ge1$ and
for $n=\infty$.
\end{theorem}

\begin{proof}
First suppose that $E$ is finite.  We use induction on
$|E\setminus A|$, simultaneously for all attaching targets $X$ and
all attaching functors $a$.  When this cardinality is zero, $j$ is an
isomorphism.  When the complement has a greatest element,
Lemma~\ref{lem:greatest-new} applies.

Otherwise the nonempty finite poset $E\setminus A$ has at least two
distinct maximal elements $x_1,x_2$: a finite poset with a unique
maximal element has that element as its greatest element.  Each $x_i$
is also maximal in $E$, since an element above a new element cannot
belong to the sieve $A$.  Put
\[
 E_1=E\setminus\{x_1\},\qquad
 E_2=E\setminus\{x_2\},\qquad
 E_{12}=E\setminus\{x_1,x_2\}.
\]
These are sieves of $E$ containing $A$; moreover
$E_1\cup E_2=E$ and $E_1\cap E_2=E_{12}$.  The restriction of $r$ to
each $E_T$, for $T=1,2,12$, still takes values in $A$ and satisfies
\eqref{eq:poset-adjunction}.  Each complement $E_T\setminus A$ is
smaller.  Define
\[
 Y_T=X\amalg_{\U{A}}\U{E_T},\qquad j_T:X\longrightarrow Y_T.
\]
By induction every $\Street j_T$ is a weak homotopy equivalence.

Corollary~\ref{cor:attachment-covers}(i), applied to the sieves
$E_1,E_2$ containing $A$ and covering $E$, identifies the canonical
maps $Y_i\to Y$ as full sieve inclusions. Their full intersection is
canonically $Y_{12}$, and they cover the objects of $Y$. These
identifications hold for the given arbitrary attaching functor $a$.

Lemma~\ref{lem:nerve-cover} now identifies the bottom pushout in the
following diagram with $\Street Y$:
\begin{equation}\label{eq:couniversal-cover-spans}
\begin{tikzcd}[column sep=large]
 \Street X \arrow[d,"\Street j_1"'] &
 \Street X \arrow[l,equal] \arrow[r,equal]
             \arrow[d,"\Street j_{12}"] &
 \Street X \arrow[d,"\Street j_2"] \\
 \Street Y_1 &
 \Street Y_{12} \arrow[l,hook] \arrow[r,hook] &
 \Street Y_2 .
\end{tikzcd}
\end{equation}
The vertical maps are the canonical maps induced by $X$, so this is
a strictly commuting map of spans.  All three are weak equivalences
by induction, and the indicated legs in both rows are
monomorphisms.  Lemma~\ref{lem:simplicial-gluing} proves that the
induced map on pushouts is a weak equivalence.  The top pushout is
$\Street X$, and its map to the bottom pushout is precisely
$\Street j$, by its restrictions to the two pieces.  This completes
the finite induction.  In particular, no compatibility between the
local cylinders on different pieces has been used.

Now let $E$ be arbitrary, and write $q=ir:E\to E$.  Since
$q^2=q$, the collection $\mathcal F$ of finite subsets $F\subset E$
with $q(F)\subset F$ is nonempty and filtered under inclusion: it
contains $\varnothing$, and finite unions remain stable.  Every
finite subset $T$ is contained in the finite stable subset
$T\cup q(T)$.  Give each $F$ its full induced order and put
$A_F=A\cap F$.  Then $A_F$ is a sieve in $F$, and stability gives
$r(F)\subset A_F$.  The restriction $r_F:F\to A_F$ is a right
adjoint retraction, since \eqref{eq:poset-adjunction} restricts to
these full subposets.  Notice that $F$ itself need not be a sieve
in $E$.

Use the restriction of $a$ to define
\[
 Y_F=X\amalg_{\U{A_F}}\U{F},\qquad j_F:X\longrightarrow Y_F.
\]
Every $j_F$ is a Thomason equivalence by the finite case.  The
inclusions $F\subset F'$ give compatible transition functors
$Y_F\to Y_{F'}$ under the same fixed $X$.  Every ordinary nerve
simplex has finite vertex support, and therefore
\[
 \OrdN E=\colim_{F\in\mathcal F}\OrdN F,
 \qquad
 \OrdN A=\colim_{F\in\mathcal F}\OrdN A_F.
\]
Applying the colimit-preserving left adjoint $c_n$ and then commuting
colimits gives
\begin{equation}\label{eq:couniversal-filtered-colimit}
 \colim_{F\in\mathcal F}Y_F
 \ \cong\ X\amalg_{\U{A}}\U{E}=Y.
\end{equation}
For clarity, the fixed copy of $X$ in this formula is essential and
legitimate: a compatible family of maps $Y_F\to Z$ consists of one
map $X\to Z$, since $\mathcal F$ is nonempty and connected, together
with compatible maps $\U{F}\to Z$ agreeing on each $\U{A_F}$.
These are exactly maps from the pushout on the right.  Under this
identification the colimit of the arrows $j_F$ is $j$.

By Lemma~\ref{lem:filtered}, $\Street$ preserves these filtered
colimits and simplicial weak homotopy equivalences are closed under
filtered colimits.  Hence $\Street j$ is a weak homotopy equivalence.
No injectivity of the transition functors $Y_F\to Y_{F'}$ is
required.  The same proof applies directly when $n=\infty$.
\end{proof}

The following unit comparison is known; see
\cite[Corollary~10.11]{AraMaltsiniotis2015},
\cite[Theorem~6.9]{Gagna2018} and
\cite[Proposition~9.2.3]{Ara2023}. We include an independent proof
using the same prism and sieve-cover method.

\begin{proposition}[The unit on every poset]\label{prop:poset-unit}
For every small poset $S$, the adjunction unit
\[
 \eta_S:\OrdN S\longrightarrow\Street\U{S}
\]
is a weak homotopy equivalence.  This assertion is natural in $S$
and holds for every finite $n\ge1$ and for $n=\infty$.
\end{proposition}

\begin{proof}
First take $S$ finite and induct on its cardinality.  If $S$ is
empty, both simplicial sets are empty and the unit is an
isomorphism.  If $S$ has a greatest element $x$, the pointwise
inequality $\id_S\le\mathrm{const}_x$ contracts $\OrdN S$.
Lemma~\ref{lem:relative-prism}, with old sieve $\varnothing$, gives
a right cylinder from $\id_{\U{S}}$ to
$c_n\OrdN(\mathrm{const}_x)$.  The latter functor factors through
$\U{\{x\}}=\Orient{0}$, the terminal strict $n$-category.
Lemma~\ref{lem:gray-nerve} therefore contracts
$\Street\U{S}$.  Both realizations are nonempty and contractible,
so $\eta_S$ is a weak homotopy equivalence.

In the remaining case choose two distinct maximal elements
$x_1,x_2$ of $S$ and put
$S_i=S\setminus\{x_i\}$ and
$S_{12}=S\setminus\{x_1,x_2\}$.  These are sieves, the two single
deletions cover $S$, and all three posets are smaller.  Their
ordinary nerves cover $\OrdN S$: a simplex belongs to a sieve
containing its last vertex.  By
Lemmas~\ref{lem:poset-pieces} and~\ref{lem:nerve-cover}, their
$\U{\,\cdot\,}$ images give a full sieve cover of $\U{S}$ and
the analogous cover of its Street nerve.  The joint zero label
identifies the intersection of the two full sieves with
$\U{S_{12}}$.  Naturality of the unit
gives a commuting map of spans
\[
\begin{tikzcd}[column sep=large]
 \OrdN S_1 \arrow[d,"\eta_{S_1}"'] &
 \OrdN S_{12} \arrow[l,hook] \arrow[r,hook]
             \arrow[d,"\eta_{S_{12}}"] &
 \OrdN S_2 \arrow[d,"\eta_{S_2}"] \\
 \Street\U{S_1} &
 \Street\U{S_{12}} \arrow[l,hook] \arrow[r,hook] &
 \Street\U{S_2} .
\end{tikzcd}
\]
The vertical maps are weak equivalences by induction; all horizontal
legs are monomorphisms.  The induced map on pushouts is a weak
equivalence by Lemma~\ref{lem:simplicial-gluing}.  Under the cover
identifications it is exactly $\eta_S$: its restriction to either
$\OrdN S_i$ agrees with that unit by naturality.  This proves the
finite case.

Finally every poset is the filtered union of its finite full
subposets.  Ordinary nerve preserves this union because each
simplex has finite support, $c_n$ preserves its colimit, and
$\Street$ preserves the resulting filtered colimit by
Lemma~\ref{lem:filtered}.  Naturality identifies $\eta_S$ with the
filtered colimit of the units on those finite subposets.
Filtered-colimit closure of simplicial weak homotopy equivalences
then proves the assertion.  This proof uses only the prism, cover,
and simplicial gluing lemmas, and in particular precedes and is
independent of the model-structure transfer.
\end{proof}
